\documentclass[11pt,a4paper]{article}
\usepackage[utf8]{inputenc}
\usepackage[T1]{fontenc}
\usepackage[catalan]{babel}
\usepackage{lmodern}
\linespread{1.22}
\usepackage[a4paper,top=2.4cm,bottom=2.2cm,left=1.5cm,right=1.5cm,headheight=1.1cm,headsep=0.5cm,footskip=1cm]{geometry}
\usepackage{amsmath,amssymb,cancel}
\usepackage{icomma}
\usepackage{graphicx}
\usepackage[table]{xcolor}
\usepackage{tikz}\usetikzlibrary{calc}
\usepackage{pgfplots}\pgfplotsset{compat=1.17}
\usepackage{fontawesome5}
\usepackage{ccicons}
\usepackage{fancyhdr}
\usepackage{enumitem}
\usepackage{booktabs,multirow}
\usepackage{tabularx}
\usepackage{needspace}
\usepackage[most]{tcolorbox}
\usepackage{titlesec}
\usepackage[hidelinks]{hyperref}
\hypersetup{pdftitle={Apunts de teoria. Unitat 10: Distribucions de probabilitat. 1r de Batxillerat CCSS},pdfsubject={Matemàtiques Aplicades a les Ciències Socials I},pdfkeywords={variable aleatòria, funció de probabilitat, esperança, variància, distribució binomial, distribució normal, campana de Gauss, normal típica, tipificació, aproximació de la binomial per la normal}}

% colors de la web (FentMates, bloc de Probabilitat i Estadística: morat; es manté el nom fmgreen dels apunts de les altres unitats)
\definecolor{fmnavy}{HTML}{1B2748}
\definecolor{fmgreen}{HTML}{684E88}
\definecolor{fmgreenl}{HTML}{8064A1}
\definecolor{fmyellow}{HTML}{FFC94D}
\definecolor{fmgold}{HTML}{E9A919}
\definecolor{fmaccent}{HTML}{FF7A50}
\definecolor{fmteal}{HTML}{2BB3A3}
\definecolor{fmdef}{HTML}{1D4F73}\definecolor{fmdefbg}{HTML}{EEF4F9}
\definecolor{fmprop}{HTML}{3B4252}\definecolor{fmpropbg}{HTML}{F3F4F7}
\definecolor{fmex}{HTML}{A3203A}\definecolor{fmexbg}{HTML}{FBF0F2}
\definecolor{fmmuted}{HTML}{5D6685}\definecolor{fmball0}{HTML}{E2483D}\definecolor{fmball1}{HTML}{2F6FD6}
\definecolor{fmsurface}{HTML}{EEF1F8}
\definecolor{fmA}{HTML}{3B82F6}\definecolor{fmB}{HTML}{F97316}\definecolor{fmR}{HTML}{E11D48}\definecolor{fmK}{HTML}{16A34A}\definecolor{fmV}{HTML}{7C3AED}

% blocs que no es poden partir: es mesuren i es reserva exactament el seu espai (amb el títol o la capçalera del requadre)
\newsavebox{\fmkb}
\newcommand{\fmmeasure}[2]{\setbox\fmkb=\vbox{\hsize=\dimexpr\linewidth-#1\relax\linewidth=\hsize #2}}
\newcommand{\fmblock}[1]{\fmmeasure{0pt}{#1}\Needspace{\dimexpr\ht\fmkb+\dp\fmkb\relax}#1}
\newcommand{\fmheadblock}[2]{\fmmeasure{0pt}{#2}\Needspace{\dimexpr\ht\fmkb+\dp\fmkb+2.2cm\relax}#1#2}
\newcommand{\fmboxneed}[1]{\fmmeasure{9mm}{#1}\Needspace{\dimexpr\ht\fmkb+\dp\fmkb+1.6cm\relax}}
\newif\iffmstack\newif\iffmfit
\newcommand{\fmfits}[1]{\par\ifdim\dimexpr\pagegoal-\pagetotal\relax<\dimexpr\ht\fmkb+\dp\fmkb+#1\relax\global\fmfitfalse\else\global\fmfittrue\fi}
\newcommand{\fmtwo}[2]{\fmmeasure{0pt}{#1}\fmfits{0pt}\iffmfit #1\else #2\fi}
\newcommand{\fmheadtwo}[3]{\fmmeasure{0pt}{#2}\fmfits{2.2cm}\iffmfit #1#2\else\Needspace{\dimexpr1.6cm+7\baselineskip\relax}#1#3\fi}
\newcommand{\fmtwoin}[2]{\iffmstack #2\else #1\fi}
\newcommand{\fmboxneedalt}[1]{\fmmeasure{9mm}{#1}\fmfits{1.6cm}\iffmfit\global\fmstackfalse\else\global\fmstacktrue\Needspace{8\baselineskip}\fi}
\newcommand{\fmheadboxalt}[2]{\fmmeasure{9mm}{#2}\fmfits{3.4cm}\iffmfit\global\fmstackfalse\else\global\fmstacktrue\Needspace{\dimexpr2.4cm+5\baselineskip\relax}\fi#1}
\newcommand{\fmheadbox}[2]{\fmmeasure{9mm}{#2}\Needspace{\dimexpr\ht\fmkb+\dp\fmkb+3.4cm\relax}#1}
\newcommand{\lt}{<}\newcommand{\gt}{>}
% fórmula en línia que no es parteix mai; si no cap en una línia, es posa centrada
% fórmula en línia; si ocupa més del 80 % de la línia, va centrada (amb la puntuació que la segueix)
\makeatletter
\newcommand{\fmm}[1]{\setbox0=\hbox{$#1$}\ifdim\wd0>0.8\linewidth\def\fmmnext{\fmmdisp{#1}}\else\def\fmmnext{\mbox{\box0}}\fi\fmmnext}
\def\fmmdisp#1{\@ifnextchar.{\fmmdispp{#1}}{\@ifnextchar,{\fmmdispc{#1}}{\@ifnextchar;{\fmmdisps{#1}}{\fmdisp{#1}}}}}
\def\fmmdispp#1.{\fmdisp{#1\,.}}\def\fmmdispc#1,{\fmdisp{#1\,,}}\def\fmmdisps#1;{\fmdisp{#1\,;}}
\makeatother
% fórmula centrada; si és més ampla que la línia, s'encongeix fins que hi cap
\newcommand{\fmdisp}[1]{\[\setbox2=\hbox{$\displaystyle #1$}\ifdim\wd2>\linewidth\resizebox{\linewidth}{!}{\box2}\else\box2\fi\]}
\DeclareUnicodeCharacter{2212}{-}
\DeclareUnicodeCharacter{2225}{\ensuremath{\parallel}}\DeclareUnicodeCharacter{22A5}{\ensuremath{\perp}}\DeclareUnicodeCharacter{2226}{\ensuremath{\nparallel}}
\DeclareUnicodeCharacter{207A}{\ensuremath{^{+}}}\DeclareUnicodeCharacter{207B}{\ensuremath{^{-}}}
\setlength{\parindent}{0pt}
\setlength{\parskip}{12pt}
\setlength{\lineskip}{5pt}\setlength{\lineskiplimit}{5pt}  % més aire al voltant de les fraccions en línia
\setlength{\emergencystretch}{3em}
\setlist{itemsep=4pt,topsep=4pt}
\renewcommand{\thesection}{10.\arabic{section}}
\titleformat{\section}{\color{fmnavy}\Large\bfseries}{\colorbox{fmdef}{\color{white}\strut\ \thesection\ }}{0.6em}{}[{\color{fmdef}\titlerule[1.2pt]}]
\titlespacing*{\section}{0pt}{10pt}{4pt}
\titleformat*{\subsection}{\color{fmgreen}\large\bfseries}
\titlespacing*{\subsection}{0pt}{6pt plus 2pt}{0pt}
\allowdisplaybreaks

% dibuix de la targeta de la unitat a la web (illustrations.js), en blanc
\newcommand{\fmcard}{\begin{tikzpicture}[x=0.062cm,y=-0.062cm,white,line width=3pt,line cap=round,line join=round]
  \begin{scope}[transparency group,opacity=0.15]
  \fill (20,60) circle[radius=5];
  \draw (20,60)--(80,26) (20,60)--(80,94) (80,26)--(130,12) (80,26)--(130,40) (80,94)--(130,80) (80,94)--(130,108);
  \fill (80,26) circle[radius=4] (80,94) circle[radius=4] (130,12) circle[radius=3.5] (130,40) circle[radius=3.5] (130,80) circle[radius=3.5] (130,108) circle[radius=3.5];
  \node[anchor=base west,inner sep=0pt,font=\fontsize{18}{18}\selectfont\itshape] at (36,36) {0,6};
  \node[anchor=base west,inner sep=0pt,font=\fontsize{18}{18}\selectfont\itshape] at (36,92) {0,4};
  \begin{scope}[rotate around={-12:(190,58)}]
    \draw[rounded corners=0.74cm] (162,30) rectangle (218,86);
    \fill (176,44) circle[radius=4.5] (190,58) circle[radius=4.5] (204,72) circle[radius=4.5];
  \end{scope}
  \end{scope}
\end{tikzpicture}}
\newcommand{\fmbrand}{\raisebox{-0.25\height}{\includegraphics[height=0.5cm]{logo-fentmates.pdf}}}

% capçalera i peu
\pagestyle{fancy}\fancyhf{}
\fancyhead[L]{\small\fmbrand}
\fancyhead[R]{\small\color{fmmuted}1r Batxillerat CCSS · Unitat 10: Distribucions de probabilitat}
\fancyfoot[L]{\scriptsize\color{fmmuted}\href{https://creativecommons.org/licenses/by-nc-sa/4.0/deed.ca}{\ccbyncsa}\ \ FentMates · fentmates.org · CC BY-NC-SA 4.0}
\fancyfoot[R]{\small\color{fmmuted}\thepage}
\renewcommand{\headrulewidth}{0.4pt}\renewcommand{\headrule}{\color{fmmuted!40}\hrule height \headrulewidth}
\fancypagestyle{primera}{\fancyhf{}\renewcommand{\headrulewidth}{0pt}
  \fancyfoot[L]{\scriptsize\color{fmmuted}\href{https://creativecommons.org/licenses/by-nc-sa/4.0/deed.ca}{\ccbyncsa}\ \ FentMates · fentmates.org · CC BY-NC-SA 4.0}
  \fancyfoot[R]{\small\color{fmmuted}\thepage}}

% caixes (mateixos colors que la web)
% quan una caixa es parteix entre dues pàgines, la vora del tall és discontínua
\tcbset{fmbase/.style={enhanced,breakable,arc=3mm,boxrule=1pt,left=3.5mm,right=3.5mm,top=2mm,bottom=2.5mm,before skip=10pt,after skip=10pt,fonttitle=\bfseries,
  frame style={fill=tcbcolframe},
  overlay first={\fill[tcbcolback] ([xshift=1pt]frame.south west) rectangle ([xshift=-1pt,yshift=1.2pt]frame.south east);
    \draw[tcbcolframe,line width=1pt,dash pattern=on 5pt off 3pt] (frame.south west)--(frame.south east);},
  overlay middle={\fill[tcbcolback] ([xshift=1pt]frame.south west) rectangle ([xshift=-1pt,yshift=1.2pt]frame.south east);
    \fill[tcbcolback] ([xshift=1pt]frame.north west) rectangle ([xshift=-1pt,yshift=-1.2pt]frame.north east);
    \draw[tcbcolframe,line width=1pt,dash pattern=on 5pt off 3pt] (frame.south west)--(frame.south east) (frame.north west)--(frame.north east);},
  overlay last={\fill[tcbcolback] ([xshift=1pt]frame.north west) rectangle ([xshift=-1pt,yshift=-1.2pt]frame.north east);
    \draw[tcbcolframe,line width=1pt,dash pattern=on 5pt off 3pt] (frame.north west)--(frame.north east);}}}
\newtcolorbox{fmdef}[1]{fmbase,colback=fmdefbg,colframe=fmdef,coltitle=white,title={#1}}
\newtcolorbox{fmprop}[1]{fmbase,colback=fmpropbg,colframe=fmprop,coltitle=white,title={#1},halign title=center}
\newtcolorbox{fmex}[1]{fmbase,before upper=\raggedright,colback=fmexbg,colframe=fmex,colbacktitle=fmexbg,coltitle=fmex,title={\faPencil*\ #1},
  titlerule=0.4pt,titlerule style={fmex!40,dashed}}
\newtcolorbox{fmwarn}[1]{fmbase,colback=fmsf!6,colframe=fmsf,coltitle=white,title={#1}}
\newtcolorbox{fmnote}[1][]{enhanced,breakable,colback=white,colframe=white,borderline west={3pt}{0pt}{fmaccent},arc=0pt,left=3mm,right=2mm,top=1mm,bottom=1mm,before skip=6pt,after skip=6pt,before upper=\raggedright,#1}
\newcommand{\fmkey}[1]{\colorbox{fmyellow!45}{\strut #1}}
\newcommand{\fmuline}[1]{\textbf{\color{fmaccent}#1}}
\definecolor{fmok}{HTML}{15803D}\definecolor{fmev}{HTML}{2563EB}\definecolor{fmsf}{HTML}{C2410C}\definecolor{fmsi}{HTML}{7C3AED}\definecolor{fmas}{HTML}{A3203A}
\newcommand{\fmbadge}[2]{\tcbox[on line,colback=#1,colframe=#1,arc=2mm,boxsep=0pt,left=4pt,right=4pt,top=1.5pt,bottom=1.5pt]{\color{white}\footnotesize\bfseries #2}}
\newcommand{\fmchip}[1]{\tcbox[on line,colback=fmsf!10,colframe=fmsf!50,arc=2mm,boxsep=1pt,left=4pt,right=4pt]{#1}}
\newcommand{\fmlevel}[2]{\tcbox[on line,colback=white,colframe=fmmuted!40,arc=2mm,boxsep=2pt,halign=center]{\begin{tabular}{c}{\scriptsize\color{fmmuted}\MakeUppercase{#1}}\\#2\end{tabular}}}
% gràfiques
\pgfplotsset{fmaxis/.style={axis lines=middle,grid=major,grid style={fmmuted!15},axis line style={fmmuted,-latex},
  tick label style={font=\scriptsize,color=fmmuted},clip=true,unbounded coords=jump,restrict y to domain=-200:200,
  every axis plot/.append style={line join=round}},
  fmcurve/.style={fmaccent,line width=1.6pt,smooth}}
\tikzset{fmasy/.style={fmmuted,dashed,line width=0.9pt},fmguide/.style={fmaccent,dotted,line width=0.8pt},
  fmdot/.style={circle,fill=fmaccent,inner sep=0pt,minimum size=5pt},fmopen/.style={circle,draw=fmaccent,fill=white,line width=1pt,inner sep=0pt,minimum size=5pt}}
\newcommand{\fmpau}[1]{\tcbox[on line,colback=fmnavy,colframe=fmnavy,arc=2mm,boxsep=0pt,left=4pt,right=4pt,top=1.5pt,bottom=1.5pt]{\color{white}\scriptsize\bfseries\MakeUppercase{#1}}\ }

\begin{document}
\thispagestyle{primera}
\begin{tikzpicture}[remember picture,overlay]
  \fill[fmgreen] (current page.north west) rectangle ([yshift=-10.1cm]current page.north east);
  \begin{scope}
    \clip (current page.north west) rectangle ([yshift=-10.1cm]current page.north east);
    \node[anchor=north west,inner sep=0pt] at ([xshift=6.6cm,yshift=-1.3cm]current page.north west) {\fmcard};
  \end{scope}
  \node[anchor=north west,inner sep=0pt] at ([xshift=1.5cm,yshift=-1.1cm]current page.north west) {\tcbox[colback=white,colframe=white,arc=2.5mm,boxsep=0pt,left=7pt,right=7pt,top=5pt,bottom=5pt]{\includegraphics[height=0.75cm]{logo-fentmates.pdf}}};
\end{tikzpicture}
\vspace*{0.55cm}
\tcbox[colback=white!15!fmgreen,colframe=white!15!fmgreen,arc=2.5mm,boxsep=0pt,left=6pt,right=6pt,top=3pt,bottom=3pt]{\color{white}\small\bfseries Apunts de teoria · 1r Batxillerat · Matemàtiques Aplicades a les Ciències Socials I}

\vspace{4pt}
{\raggedright\color{white}\fontsize{25}{30}\selectfont\bfseries{\color{fmyellow}Unitat 10} · Distribucions de probabilitat\par}
\vspace{6pt}
\begin{minipage}{\linewidth}\color{white!88!fmgreen}\small Teoria i exemples resolts pas a pas: variables aleatòries discretes i contínues, la funció de probabilitat i la funció de distribució, la mitjana i la desviació típica d'una variable aleatòria, la distribució binomial (amb la fórmula i amb la taula), la funció de densitat i la distribució uniforme, la distribució normal, la normal típica \fmm{N(0,1)} i la seua taula, la tipificació i l'aproximació de la distribució binomial a la normal.\end{minipage}\par
% l'índex comença sempre a la mateixa altura, per davall de la franja (com en les unitats 5-10)
\vspace*{\dimexpr 8.6cm-\pagetotal\relax}

\begin{tcolorbox}[colback=fmsurface,colframe=fmsurface,arc=3mm,left=4mm,right=4mm,top=2mm,bottom=2mm]
\small\textbf{Índex}\par\smallskip
\begin{minipage}[t]{0.49\linewidth}\small \mbox{\hyperlink{sec1}{\textbf{\color{fmgreen}10.1}\ \ Variables aleatòries}}\\[3pt]
\mbox{\hyperlink{sec2}{\textbf{\color{fmgreen}10.2}\ \ Variables aleatòries discretes}}\\[3pt]
\mbox{\hyperlink{sec3}{\textbf{\color{fmgreen}10.3}\ \ La distribució binomial}}\\[3pt]
\mbox{\hyperlink{sec4}{\textbf{\color{fmgreen}10.4}\ \ La taula de la binomial}}\\[3pt]
\mbox{\hyperlink{sec5}{\textbf{\color{fmgreen}10.5}\ \ Variables aleatòries contínues}}\\[3pt]\end{minipage}\hfill\begin{minipage}[t]{0.49\linewidth}\small \mbox{\hyperlink{sec6}{\textbf{\color{fmgreen}10.6}\ \ La distribució normal}}\\[3pt]
\mbox{\hyperlink{sec7}{\textbf{\color{fmgreen}10.7}\ \ La normal típica i la seua taula}}\\[3pt]
\mbox{\hyperlink{sec8}{\textbf{\color{fmgreen}10.8}\ \ Tipificació. Problemes}}\\[3pt]
\hyperlink{sec9}{\textbf{\color{fmgreen}10.9}\ \ Aproximació de la binomial a la normal}\\[3pt]\mbox{\hyperlink{annex}{\textbf{\color{fmgreen}Annex}\ \ Taules}}\end{minipage}
\end{tcolorbox}

\Needspace{8\baselineskip}
\section{Variables aleatòries. Discretes i contínues}\hypertarget{sec1}{}

A la unitat 9 hem calculat probabilitats de \textbf{successos}. Ara farem un pas més: en molts experiments el que ens interessa és un \textbf{nombre} (quantes cares ixen, quants clients arriben, quant pesa una peça\dots{}). Si a cada resultat li associem un nombre, podem estudiar l'experiment amb les eines de l'estadística: taules, gràfics, mitjana i desviació típica. Una \textbf{distribució de probabilitat} és, per dir-ho així, una \emph{distribució de freqüències} «teòrica»: la que obtindríem si repetírem l'experiment infinites vegades.

\Needspace{5\baselineskip}
\begin{fmdef}{Variable aleatòria}
Una \textbf{variable aleatòria} \fmm{X} és una funció que assigna un nombre real a cada resultat d'un experiment aleatori. S'escriu amb majúscula (\fmm{X}, \fmm{Y}\dots{}) i els valors que pren, amb minúscula (\fmm{x_1, x_2,\dots}).

Per exemple, en llançar tres monedes, \fmm{X=}«nombre de cares» assigna \fmm{X(CCC)=3}, \fmm{X(CCX)=2}, \dots{}, \fmm{X(XXX)=0}: pren els valors 0, 1, 2 i 3.

\par\noindent \begin{minipage}[t]{0.47\linewidth}\vspace{0pt}\raggedright \textbf{Discreta}: pren un nombre finit (o numerable) de valors aïllats, que es poden comptar.

\begin{itemize}
\item Nombre de cares en 10 llançaments.
\item Nombre de clients que entren en una hora.
\item Nombre de peces defectuoses d'una caixa.
\end{itemize}

\begin{fmnote}Es descriu amb una \textbf{funció de probabilitat} (taula i diagrama de barres).\end{fmnote}\end{minipage}\hfill\begin{minipage}[t]{0.46\linewidth}\vspace{0pt}\raggedright \textbf{Contínua}: pot prendre qualsevol valor d'un interval (és una mesura).

\begin{itemize}
\item L'altura o el pes d'una persona.
\item El temps d'espera a la parada de l'autobús.
\item La durada d'una bateria.
\end{itemize}

\begin{fmnote}Es descriu amb una \textbf{funció de densitat} (una corba; la probabilitat és una àrea).\end{fmnote}\end{minipage}\par\medskip
\end{fmdef}

\Needspace{5\baselineskip}
\begin{fmprop}{Estadística i probabilitat: el para\lgem{}elisme}
\begin{center}\renewcommand{\arraystretch}{1.5}\setlength{\parskip}{0pt}\small\begin{tabularx}{\linewidth}{>{\raggedright\arraybackslash}p{5.2cm}>{\raggedright\arraybackslash}X}\toprule
\textbf{Estadística (unitat 8 i anteriors)} & \textbf{Probabilitat (aquesta unitat)}\\\midrule
Variable estadística & Variable aleatòria\\[3pt]\end{tabularx}\par\nointerlineskip\vskip0pt
\small\begin{tabularx}{\linewidth}{>{\raggedright\arraybackslash}p{5.2cm}>{\raggedright\arraybackslash}X}Freqüència relativa $f_i$ & Probabilitat $p_i=P(X=x_i)$\\[3pt]\end{tabularx}\par\nointerlineskip\vskip0pt
\small\begin{tabularx}{\linewidth}{>{\raggedright\arraybackslash}p{5.2cm}>{\raggedright\arraybackslash}X}Freqüència relativa acumulada $F_i$ & Funció de distribució $F(x)=P(X\leq x)$\\[3pt]\end{tabularx}\par\nointerlineskip\vskip0pt
\small\begin{tabularx}{\linewidth}{>{\raggedright\arraybackslash}p{5.2cm}>{\raggedright\arraybackslash}X}Diagrama de barres / histograma & Diagrama de barres / funció de densitat\\[3pt]\end{tabularx}\par\nointerlineskip\vskip0pt
\small\begin{tabularx}{\linewidth}{>{\raggedright\arraybackslash}p{5.2cm}>{\raggedright\arraybackslash}X}Mitjana $\bar x=\sum x_i f_i$ & Mitjana o esperança $\mu=\sum x_i p_i$\\[3pt]\end{tabularx}\par\nointerlineskip\vskip0pt
\small\begin{tabularx}{\linewidth}{>{\raggedright\arraybackslash}p{5.2cm}>{\raggedright\arraybackslash}X}Desviació típica $s$ & Desviació típica $\sigma$\\[3pt]\bottomrule\end{tabularx}\end{center}

\begin{fmnote}Llei dels grans nombres: si repetim l'experiment moltes vegades, les freqüències relatives s'acosten a les probabilitats i la mitjana dels resultats s'acosta a \fmm{\mu}.\end{fmnote}
\end{fmprop}

\enlargethispage{1cm}\Needspace{4cm}
\section{Variables aleatòries discretes. Funció de probabilitat, mitjana i desviació típica}\hypertarget{sec2}{}
\begin{fmdef}{Funció de probabilitat}
La \textbf{funció de probabilitat} (o \textbf{funció de massa}) d'una variable discreta assigna a cada valor \fmm{x_i} la seua probabilitat \fmm{p_i=P(X=x_i)}. Normalment es dona amb una taula i es representa amb un \textbf{diagrama de barres}. Sempre es compleix:

\fmdisp{0\leq p_i\leq1\qquad\qquad p_1+p_2+\dots+p_n=\sum p_i=1}

La \textbf{funció de distribució} dona la probabilitat acumulada: \fmm{F(x)=P(X\leq x)}. Es calcula sumant les probabilitats dels valors menors o iguals que \fmm{x}; és creixent i val 1 a partir de l'últim valor.
\end{fmdef}

\Needspace{5\baselineskip}
\begin{fmdef}{Paràmetres: mitjana, variància i desviació típica}

\fmdisp{\mu=E(X)=\sum x_i\,p_i\qquad\qquad\sigma^2=\sum (x_i-\mu)^2\,p_i=\sum x_i^2\,p_i-\mu^2\qquad\qquad\sigma=\sqrt{\sigma^2}}

\begin{itemize}
\item La \textbf{mitjana} o \textbf{esperança} \fmm{\mu} és el valor que esperem obtindre «de mitjana» si repetim l'experiment moltes vegades. No cal que siga un dels valors de \fmm{X} (de mitjana, en llançar tres monedes ixen 1,5 cares).
\item La \textbf{desviació típica} \fmm{\sigma} mesura com de dispersos estan els valors al voltant de la mitjana. La segona fórmula de \fmm{\sigma^2} és més còmoda per a fer els càlculs.
\item En un \textbf{joc}, si \fmm{X} és el guany, \fmm{\mu} és el guany mitjà per partida. El joc és \textbf{just} (o equitatiu) si \fmm{\mu=0}; si \fmm{\mu\lt0}, a la llarga el jugador perd.
\end{itemize}

\begin{fmnote}Per a calcular-ho tot de colp, afig a la taula les files \fmm{x_i\,p_i} i \fmm{x_i^2\,p_i} i suma-les (com fèiem amb les freqüències).\end{fmnote}
\end{fmdef}

\Needspace{5\baselineskip}
\begin{fmex}{Exemples resolts}
\begin{enumerate}[label=\textbf{\arabic*.},leftmargin=*]
\item Llancem tres monedes i \fmm{X=}«nombre de cares». Troba la funció de probabilitat, la funció de distribució, la mitjana i la desviació típica.

Dels 8 resultats equiprobables (CCC, CCX, \dots{}, XXX), n'hi ha 1, 3, 3 i 1 amb 0, 1, 2 i 3 cares. Fem la taula amb les files per als càlculs:

\begin{center}\renewcommand{\arraystretch}{1.5}\begin{tabular}{cccccc}\toprule
\textbf{$x_i$} & 0 & 1 & 2 & 3 & \textbf{$\Sigma$}\\
\textbf{$p_i$} & $\frac18$ & $\frac38$ & $\frac38$ & $\frac18$ & 1\\
\textbf{$F(x_i)$} & $\frac18$ & $\frac48$ & $\frac78$ & 1 & \\
\textbf{$x_i\,p_i$} & 0 & $\frac38$ & $\frac68$ & $\frac38$ & $\frac{12}8$\\
\textbf{$x_i^2\,p_i$} & 0 & $\frac38$ & $\frac{12}8$ & $\frac98$ & $\frac{24}8$\\
\bottomrule\end{tabular}\end{center}

\fmm{\mu=\frac{12}{8}=1{,}5} cares

\fmm{\sigma^2=\frac{24}{8}-1{,}5^2=3-2{,}25=0{,}75\Rightarrow\sigma=\sqrt{0{,}75}\approx0{,}866}

Per exemple, \fmm{P(X\leq1)=F(1)=\frac48=0{,}5} i \fmm{P(X\geq2)=1-F(1)=0{,}5}.
\item Una variable aleatòria pren els valors 1, 2, 3 i 4 amb probabilitats \fmm{0{,}1}, \fmm{0{,}3}, \fmm{k} i \fmm{2k}. Calcula \fmm{k}, la mitjana i \fmm{P(X\gt2)}.

\fmm{0{,}1+0{,}3+k+2k=1\Rightarrow3k=0{,}6\Rightarrow k=0{,}2}

\fmm{\mu=1\cdot0{,}1+2\cdot0{,}3+3\cdot0{,}2+4\cdot0{,}4=0{,}1+0{,}6+0{,}6+1{,}6=2{,}9}

\fmm{P(X\gt2)=P(X=3)+P(X=4)=0{,}2+0{,}4=0{,}6}
\item En una fira, per jugar es paga 2 €. Es llancen dos daus: si la suma és 12 es reben 30 €; si és 7, es reben 6 €; en cas contrari no es rep res. És un joc just?

\fmm{X=}«guany net». \fmm{P(12)=\frac1{36}}, \fmm{P(7)=\frac6{36}}, \fmm{P(\text{altre})=\frac{29}{36}}. Els guanys nets són \fmm{28}, \fmm{4} i \fmm{-2} €.

\fmm{\mu=28\cdot\frac1{36}+4\cdot\frac6{36}-2\cdot\frac{29}{36}=\frac{28+24-58}{36}=-\frac{6}{36}\approx-0{,}17} €

No és just: de mitjana el jugador perd uns 17 cèntims per partida (és el benefici del firaire).
\end{enumerate}
\end{fmex}

\Needspace{5.7cm}
\section{La distribució binomial}\hypertarget{sec3}{}
\begin{fmdef}{Quan és binomial?}
Una variable \fmm{X} segueix una \textbf{distribució binomial} \fmm{B(n,p)} quan compta el \textbf{nombre d'èxits} en un experiment que compleix:

\begin{enumerate}
\item Es repeteix una mateixa prova \fmm{n} vegades.
\item Cada prova només té dos resultats: \textbf{èxit} (\fmm{A}) o \textbf{fracàs} (\fmm{\overline A}).
\item La probabilitat d'èxit, \fmm{p=P(A)}, és la mateixa en totes les proves (la de fracàs és \fmm{q=1-p}).
\item Les proves són \textbf{independents} (per exemple, extraccions \textbf{amb} reemplaçament, o d'una població molt gran).
\end{enumerate}

\begin{fmnote}Exemples: nombre de cares en 10 llançaments, nombre de peces defectuoses en una mostra de 20 (si el 3 \% ho són), nombre d'encerts en un test contestat a l'atzar\dots{} Compte: les extraccions \textbf{sense} reemplaçament d'una urna xicoteta \textbf{no} són binomials (la \fmm{p} canvia).\end{fmnote}
\end{fmdef}

\Needspace{5\baselineskip}
\begin{fmex}{D'on ix la fórmula?}
Un jugador de bàsquet encistella un tir lliure amb probabilitat \fmm{p=0{,}7}. Llança 3 tirs (independents). Quina és la probabilitat que n'encerte exactament 2?

Els resultats amb 2 encerts (E) i 1 fallada (F) són \fmm{EEF}, \fmm{EFE} i \fmm{FEE}. Per la independència, cada un té probabilitat \fmm{0{,}7\cdot0{,}7\cdot0{,}3=0{,}7^2\cdot0{,}3^1=0{,}147}.

Quants n'hi ha? Tantes com maneres de triar en quines 2 de les 3 proves hi ha l'encert: \fmm{\binom32=3} (combinacions, unitat 9).

\fmm{P(X=2)=\binom32\cdot0{,}7^2\cdot0{,}3^1=3\cdot0{,}147=0{,}441}
\end{fmex}

\Needspace{5\baselineskip}
\begin{fmprop}{Fórmula de la binomial}

\fmdisp{P(X=k)=\binom{n}{k}\,p^k\,q^{\,n-k},\qquad k=0,1,\dots,n\qquad\qquad\binom nk=\frac{n!}{k!\,(n-k)!}}

El \textbf{nombre combinatori} \fmm{\binom nk} («\fmm{n} sobre \fmm{k}») compta de quantes maneres es poden co\lgem{}ocar els \fmm{k} èxits entre les \fmm{n} proves; \fmm{p^kq^{n-k}} és la probabilitat de cada una d'aquestes maneres. Recorda: \fmm{\binom n0=\binom nn=1}, \fmm{\binom n1=n}, \fmm{\binom nk=\binom n{n-k}}. A la calculadora: tecla \textbf{nCr}.

\fmdisp{\mu=n\,p\qquad\qquad\sigma=\sqrt{n\,p\,q}}

Per a probabilitats acumulades sumem: \fmm{P(X\leq k)=P(X=0)+\dots+P(X=k)}. I quan «almenys un» o «més de\dots{}» tinga molts termes, fem servir el contrari: \fmm{P(X\geq1)=1-P(X=0)}.
\end{fmprop}

\Needspace{5\baselineskip}
\begin{fmex}{Exemples resolts}
\begin{enumerate}[label=\textbf{\arabic*.},leftmargin=*]
\item Un test té 10 preguntes amb 4 opcions cada una. Un alumne contesta totes a l'atzar. Probabilitat que n'encerte exactament 3, i que n'encerte almenys 1.

\fmm{X=}«nombre d'encerts» és \fmm{B(10;\,0{,}25)}: 10 proves independents amb \fmm{p=\frac14} d'encert.

\fmm{P(X=3)=\binom{10}{3}\,0{,}25^3\,0{,}75^7=120\cdot0{,}015625\cdot0{,}1335\approx0{,}2503}

\fmm{P(X\geq1)=1-P(X=0)=1-0{,}75^{10}\approx1-0{,}0563=0{,}9437}

Mitjana: \fmm{\mu=10\cdot0{,}25=2{,}5} encerts; \fmm{\sigma=\sqrt{10\cdot0{,}25\cdot0{,}75}\approx1{,}37}.
\item Un 48 \% dels habitants d'una ciutat són homes. Triem 4 persones amb reemplaçament i \fmm{X} compta els homes. Quina distribució segueix \fmm{X}? Calcula \fmm{P(X=2)}. \fmpau{PAU Aragó 2024}

\fmm{X\sim B(4;\,0{,}48)}; \fmm{\ P(X=2)=\binom42\,0{,}48^2\,0{,}52^2=6\cdot0{,}2304\cdot0{,}2704\approx0{,}3738}
\end{enumerate}
\end{fmex}

\Needspace{5.7cm}
\section{La taula de la distribució binomial}\hypertarget{sec4}{}
\begin{fmdef}{Com es llig la taula}
Per a no haver de fer tants càlculs, la probabilitat de la binomial es pot llegir en una taula (és la mateixa que s'adjunta a l'examen de la PAU valenciana, que faràs l'any que ve). Dona directament \fmm{P(X=x)} per a \fmm{n=2,\dots,9} i els valors de \fmm{p} de la capçalera. Busquem el bloc de la \fmm{n}, la fila de la \fmm{x} i la columna de la \fmm{p}.

\begin{itemize}
\item \textbf{Probabilitats acumulades}: sumem diverses caselles de la mateixa columna, o fem el contrari.
\item \textbf{Si \fmm{p\gt0{,}5}} (no és a la taula): comptem els fracassos. Si \fmm{X\sim B(n,p)}, aleshores \fmm{Y=n-X\sim B(n,q)} amb \fmm{q=1-p\lt0{,}5}, i \fmm{P(X=k)=P(Y=n-k)}.
\item \textbf{Si \fmm{n} o \fmm{p} no són a la taula} (\fmm{n\geq10}, \fmm{p=0{,}62}\dots{}): fem servir la fórmula; i si \fmm{n} és gran, l'aproximació per la normal (apartat 10.9).
\end{itemize}
\end{fmdef}

\begin{fmnote}\textbf{La taula.} La taula completa, la mateixa que s'adjunta a l'examen de la PAU, és a l'annex (pàgina \pageref{tbin}).\end{fmnote}

\Needspace{5\baselineskip}
\begin{fmex}{Exemples resolts amb la taula}
\begin{enumerate}[label=\textbf{\arabic*.},leftmargin=*]
\item Un 20 \% dels clients d'una botiga paga amb el mòbil. Triem 8 clients. Probabilitat que com a molt 2 paguen amb el mòbil.

\fmm{X\sim B(8;\,0{,}20)}. Columna \fmm{p=0{,}20}, bloc \fmm{n=8}:

\fmm{P(X\leq2)=P(0)+P(1)+P(2)=0{,}1678+0{,}3355+0{,}2936=0{,}7969}
\item El 70 \% dels alumnes aprova. Triem 6 alumnes. Probabilitat que n'aproven exactament 5, i que n'aproven almenys 5.

\fmm{p=0{,}7} no és a la taula. Comptem els suspesos: \fmm{Y\sim B(6;\,0{,}30)}, i «aproven 5» és «suspén 1».

\fmm{P(X=5)=P(Y=1)=0{,}3025}

\fmm{P(X\geq5)=P(Y\leq1)=P(Y=0)+P(Y=1)=0{,}1176+0{,}3025=0{,}4201}
\item El 87 \% dels joves compra per Internet. En una enquesta a 10 joves, probabilitat que exactament 2 no hi compren. \fmpau{PAU Castella i Lleó 2026 (adaptat)}

\fmm{X=}«no compren» \fmm{\sim B(10;\,0{,}13)}: ni \fmm{n=10} ni \fmm{p=0{,}13} són a la taula. Fórmula:

\fmm{P(X=2)=\binom{10}2\,0{,}13^2\,0{,}87^8=45\cdot0{,}0169\cdot0{,}3282\approx0{,}2496}
\end{enumerate}
\end{fmex}

\Needspace{8\baselineskip}
\section{Variables aleatòries contínues. Funció de densitat}\hypertarget{sec5}{}

Si una variable pot prendre qualsevol valor d'un interval, no podem fer una taula amb tots els valors. Passa com en estadística amb les dades agrupades en intervals: l'\textbf{histograma} dona les freqüències amb \textbf{àrees}. Si fem els intervals cada vegada més estrets, l'histograma s'acosta a una corba: la \textbf{funció de densitat}.

\fmboxneedalt{\par\noindent \begin{minipage}[t]{0.48\linewidth}\vspace{0pt}\raggedright Una variable contínua \fmm{X} es descriu amb una \textbf{funció de densitat} \fmm{f(x)} que compleix:

\begin{itemize}
\item \fmm{f(x)\geq0} per a tot \fmm{x} (la corba no baixa de l'eix).
\item L'\textbf{àrea total} entre la corba i l'eix \fmm{OX} és 1.
\end{itemize}

\fmdisp{P(a\leq X\leq b)=\text{àrea davall de }f\text{ entre }a\text{ i }b}

\begin{fmnote}La probabilitat d'un valor concret és 0 (l'àrea d'un segment és 0): \fmm{P(X=a)=0}. Per tant, en una variable contínua \textbf{és igual} posar \fmm{\lt} o \fmm{\leq}: \fmm{P(X\lt a)=P(X\leq a)}. En una variable discreta, en canvi, \textbf{no} és el mateix!\end{fmnote}\end{minipage}\hfill\begin{minipage}[t]{0.48\linewidth}\vspace{0pt}\raggedright \par\medskip\begin{minipage}{\linewidth}{\centering \includegraphics[width=\linewidth,height=4.5cm,keepaspectratio]{figs/8af8f8137419.pdf}\par}\par\smallskip{\centering\footnotesize La probabilitat és l'àrea ombrejada\par}\end{minipage}\par\medskip\end{minipage}\par\medskip}
\begin{fmdef}{Funció de densitat}
\fmtwoin{\par\noindent \begin{minipage}[t]{0.48\linewidth}\vspace{0pt}\raggedright Una variable contínua \fmm{X} es descriu amb una \textbf{funció de densitat} \fmm{f(x)} que compleix:

\begin{itemize}
\item \fmm{f(x)\geq0} per a tot \fmm{x} (la corba no baixa de l'eix).
\item L'\textbf{àrea total} entre la corba i l'eix \fmm{OX} és 1.
\end{itemize}

\fmdisp{P(a\leq X\leq b)=\text{àrea davall de }f\text{ entre }a\text{ i }b}

\begin{fmnote}La probabilitat d'un valor concret és 0 (l'àrea d'un segment és 0): \fmm{P(X=a)=0}. Per tant, en una variable contínua \textbf{és igual} posar \fmm{\lt} o \fmm{\leq}: \fmm{P(X\lt a)=P(X\leq a)}. En una variable discreta, en canvi, \textbf{no} és el mateix!\end{fmnote}\end{minipage}\hfill\begin{minipage}[t]{0.48\linewidth}\vspace{0pt}\raggedright \par\medskip\begin{minipage}{\linewidth}{\centering \includegraphics[width=\linewidth,height=4.5cm,keepaspectratio]{figs/8af8f8137419.pdf}\par}\par\smallskip{\centering\footnotesize La probabilitat és l'àrea ombrejada\par}\end{minipage}\par\medskip\end{minipage}\par\medskip}{

Una variable contínua \fmm{X} es descriu amb una \textbf{funció de densitat} \fmm{f(x)} que compleix:

\begin{itemize}
\item \fmm{f(x)\geq0} per a tot \fmm{x} (la corba no baixa de l'eix).
\item L'\textbf{àrea total} entre la corba i l'eix \fmm{OX} és 1.
\end{itemize}

\fmdisp{P(a\leq X\leq b)=\text{àrea davall de }f\text{ entre }a\text{ i }b}

\begin{fmnote}La probabilitat d'un valor concret és 0 (l'àrea d'un segment és 0): \fmm{P(X=a)=0}. Per tant, en una variable contínua \textbf{és igual} posar \fmm{\lt} o \fmm{\leq}: \fmm{P(X\lt a)=P(X\leq a)}. En una variable discreta, en canvi, \textbf{no} és el mateix!\end{fmnote}

\begin{center}\begin{minipage}{0.62\linewidth}\centering \par\medskip\begin{minipage}{\linewidth}{\centering \includegraphics[width=\linewidth,height=4.5cm,keepaspectratio]{figs/8af8f8137419.pdf}\par}\par\smallskip{\centering\footnotesize La probabilitat és l'àrea ombrejada\par}\end{minipage}\par\medskip\end{minipage}\end{center}
}

La \textbf{funció de distribució} és, com en les discretes, \fmm{F(x)=P(X\leq x)}: l'àrea a l'esquerra de \fmm{x}. Amb ella, \fmm{P(a\leq X\leq b)=F(b)-F(a)}. Les variables contínues també tenen mitjana \fmm{\mu} i desviació típica \fmm{\sigma}, que s'interpreten igual que en les discretes. En aquesta unitat les àrees seran de figures senzilles (rectangles, triangles, trapezis) o les llegirem en una taula.
\end{fmdef}\global\fmstackfalse

\Needspace{5\baselineskip}
\begin{fmprop}{La distribució uniforme $U(a,b)$}
Una variable és \textbf{uniforme} en \fmm{[a,b]} si tots els subintervals de la mateixa longitud són igual de probables. La densitat és constant: un rectangle de base \fmm{b-a} i altura \fmm{\frac1{b-a}} (perquè l'àrea siga 1).

\fmdisp{f(x)=\begin{cases}\dfrac{1}{b-a} & \text{si } a\leq x\leq b\\ 0 & en cas contrari\end{cases}\qquad P(c\leq X\leq d)=\frac{d-c}{b-a}\qquad\mu=\frac{a+b}{2}}

\end{fmprop}

\Needspace{5\baselineskip}
\begin{fmex}{Exemples resolts}
\begin{enumerate}[label=\textbf{\arabic*.},leftmargin=*]
\item Un autobús passa cada 15 minuts. Una persona arriba a la parada sense mirar l'horari, de manera que el temps d'espera \fmm{X} (en minuts) és uniforme en \fmm{[0,15]}. Probabilitat que espere menys de 4 minuts, i que espere entre 5 i 10 minuts. Temps mitjà d'espera?

Densitat: \fmm{f(x)=\frac1{15}} en \fmm{[0,15]}. Cada probabilitat és l'àrea d'un rectangle (base \fmm{\times} altura):

\fmm{P(X\lt4)=4\cdot\frac1{15}=\frac4{15}\approx0{,}267\qquad P(5\leq X\leq10)=5\cdot\frac1{15}=\frac13\approx0{,}333}

\fmm{\mu=\frac{0+15}{2}=7{,}5} minuts.
\item Comprova que \fmm{f(x)=\frac x2} per a \fmm{0\leq x\leq2} (i 0 fora) és una funció de densitat i calcula \fmm{P(X\leq1)} i \fmm{P(1\leq X\leq1{,}5)}.

\fmm{f(x)\geq0} i la regió és un triangle de base 2 i altura \fmm{f(2)=1}: àrea \fmm{\frac{2\cdot1}{2}=1}.  És una densitat.

\fmm{P(X\leq1)}: triangle de base 1 i altura \fmm{f(1)=0{,}5}: \fmm{\frac{1\cdot0{,}5}{2}=0{,}25}.

\fmm{P(1\leq X\leq1{,}5)}: trapezi de bases \fmm{f(1)=0{,}5} i \fmm{f(1{,}5)=0{,}75} i altura \fmm{0{,}5}: \fmm{\frac{0{,}5+0{,}75}{2}\cdot0{,}5=0{,}3125}.
\end{enumerate}
\end{fmex}

\Needspace{5.7cm}
\section{La distribució normal $N(\mu,\sigma)$}\hypertarget{sec6}{}
\begin{fmdef}{La campana de Gauss}
Moltes variables reals (altures, pesos, notes, temps, errors de mesura\dots{}) segueixen una \textbf{distribució normal} de mitjana \fmm{\mu} i desviació típica \fmm{\sigma}, \fmm{X\sim N(\mu,\sigma)}. La seua densitat és la \textbf{campana de Gauss}:

\fmdisp{f(x)=\frac{1}{\sigma\sqrt{2\pi}}\,e^{-\frac{(x-\mu)^2}{2\sigma^2}}}

\begin{itemize}
\item És \textbf{simètrica} respecte de \fmm{x=\mu}, on té el màxim: \fmm{P(X\leq\mu)=P(X\geq\mu)=0{,}5}.
\item \fmm{\mu} indica on està centrada; \fmm{\sigma}, com d'ampla és: com més gran \fmm{\sigma}, més baixa i ampla.
\item Regla 68--95--99,7: entre \fmm{\mu\pm\sigma} hi ha aproximadament el 68 \% dels valors; entre \fmm{\mu\pm2\sigma}, el 95 \%; entre \fmm{\mu\pm3\sigma}, el 99,7 \%.
\end{itemize}

\begin{fmnote}L'àrea davall de la campana no es pot calcular amb figures geomètriques: les probabilitats d'una normal es llegeixen en la \textbf{taula de la normal típica} (apartats 10.7 i 10.8).\end{fmnote}
\end{fmdef}

\Needspace{5.7cm}
\section{La normal típica $N(0,1)$ i la seua taula}\hypertarget{sec7}{}
\begin{fmdef}{Com es llig la taula}
La \textbf{normal típica} és \fmm{Z\sim N(0,1)}. La taula dona \fmm{F(z)=P(Z\leq z)}, l'àrea a l'\textbf{esquerra} de \fmm{z}, per a \fmm{z} entre 0,00 i 3,99. La fila dona les unitats i les dècimes de \fmm{z} i la columna, les centèsimes. Per exemple, \fmm{P(Z\leq1{,}25)}: fila 1,2 i columna 0,05 \fmm{\to0{,}8944}.

Amb la simetria i el contrari es calcula tot:

\par\noindent \begin{minipage}[t]{0.31\linewidth}\centering {\centering \includegraphics[width=\linewidth,height=3.4cm,keepaspectratio]{figs/377a3769f3f3.pdf}\par}\par\smallskip{\centering\footnotesize \fmm{P(Z\leq z)}: taula\par}\end{minipage}\hfill\begin{minipage}[t]{0.31\linewidth}\centering {\centering \includegraphics[width=\linewidth,height=3.4cm,keepaspectratio]{figs/ab3e47b48d55.pdf}\par}\par\smallskip{\centering\footnotesize \fmm{P(Z\geq z)=1-P(Z\leq z)}\par}\end{minipage}\hfill\begin{minipage}[t]{0.31\linewidth}\centering {\centering \includegraphics[width=\linewidth,height=3.4cm,keepaspectratio]{figs/20c6929f1e49.pdf}\par}\par\smallskip{\centering\footnotesize \fmm{P(Z\leq -z)=1-P(Z\leq z)}\par}\end{minipage}\par\medskip
\par\noindent \begin{minipage}[t]{0.31\linewidth}\centering {\centering \includegraphics[width=\linewidth,height=3.4cm,keepaspectratio]{figs/2a4c02e15a8e.pdf}\par}\par\smallskip{\centering\footnotesize \fmm{P(Z\geq -z)=P(Z\leq z)}\par}\end{minipage}\hfill\begin{minipage}[t]{0.31\linewidth}\centering {\centering \includegraphics[width=\linewidth,height=3.4cm,keepaspectratio]{figs/edc96f2418bd.pdf}\par}\par\smallskip{\centering\footnotesize \fmm{P(a\leq Z\leq b)=P(Z\leq b)-P(Z\leq a)}\par}\end{minipage}\hfill\begin{minipage}[t]{0.31\linewidth}\centering {\centering \includegraphics[width=\linewidth,height=3.4cm,keepaspectratio]{figs/7f428a401cb8.pdf}\par}\par\smallskip{\centering\footnotesize \fmm{P(-z\leq Z\leq z)=2P(Z\leq z)-1}\par}\end{minipage}\par\medskip

\begin{fmnote}\fmuline{Problema invers}: si ens donen la probabilitat i volem \fmm{z}, la busquem \textbf{dins} de la taula i llegim la fila i la columna (si no hi és exacta, prenem la més pròxima). Si la probabilitat és menor que 0,5, \fmm{z} és negatiu: busquem \fmm{1-p} i canviem el signe.\end{fmnote}
\end{fmdef}

\begin{fmnote}\textbf{La taula.} La taula completa, la mateixa que s'adjunta a l'examen de la PAU, és a l'annex (pàgina \pageref{tnorm}).\end{fmnote}

\Needspace{5\baselineskip}
\begin{fmex}{Exemples resolts}
\begin{enumerate}[label=\textbf{\arabic*.},leftmargin=*]
\item Si \fmm{Z\sim N(0,1)}, calcula \fmm{P(Z\gt1{,}25)}, \fmm{P(Z\lt-0{,}84)} i \fmm{P(-1\lt Z\lt2)}.

\fmm{P(Z\gt1{,}25)=1-P(Z\leq1{,}25)=1-0{,}8944=0{,}1056}

Per simetria: \fmm{P(Z\lt-0{,}84)=P(Z\gt0{,}84)=1-0{,}7995=0{,}2005}

\fmm{P(-1\lt Z\lt2)=P(Z\lt2)-P(Z\lt-1)=0{,}9772-(1-0{,}8413)=0{,}8185}
\item Troba \fmm{k} tal que \fmm{P(Z\leq k)=0{,}9} i \fmm{c} tal que \fmm{P(Z\geq c)=0{,}95}.

Dins de la taula, el valor més pròxim a 0,9 és 0,8997 (fila 1,2, columna 0,08): \fmm{k\approx1{,}28}.

\fmm{P(Z\geq c)=0{,}95\gt0{,}5}: \fmm{c} és negatiu i, per simetria, \fmm{P(Z\leq-c)=0{,}95}. A la taula, 0,9495 (1,64) i 0,9505 (1,65) estan igual de prop: \fmm{-c\approx1{,}645\Rightarrow c\approx-1{,}645}.
\end{enumerate}
\end{fmex}

\Needspace{5.7cm}
\section{Tipificació. Problemes amb la normal}\hypertarget{sec8}{}
\begin{fmprop}{Tipificació}
Si \fmm{X\sim N(\mu,\sigma)}, la variable \fmm{Z=\dfrac{X-\mu}{\sigma}} segueix una \fmm{N(0,1)}. Per tant:

\fmdisp{P(X\leq a)=P\!\left(Z\leq\frac{a-\mu}{\sigma}\right)\qquad\qquad P(a\leq X\leq b)=P\!\left(\frac{a-\mu}{\sigma}\leq Z\leq\frac{b-\mu}{\sigma}\right)}

El valor \fmm{z=\frac{a-\mu}{\sigma}} diu a quantes desviacions típiques està \fmm{a} de la mitjana. Arredonim \fmm{z} a dues decimals per a poder fer servir la taula (com a la PAU).

\begin{fmnote}\fmuline{Tipus de problemes}: (1) \textbf{directes}: ens donen \fmm{a} i volem la probabilitat; (2) \textbf{«quantes persones\dots{}»}: multipliquem la probabilitat pel total; (3) \textbf{inversos}: ens donen la probabilitat i volem el valor \fmm{a} (\fmm{a=\mu+z\sigma}), o bé \fmm{\mu} o \fmm{\sigma}.\end{fmnote}
\end{fmprop}

\Needspace{5\baselineskip}
\begin{fmex}{Exemples resolts}
\begin{enumerate}[label=\textbf{\arabic*.},leftmargin=*]
\item El coeficient inte\lgem{}ectual d'una població de 1200 persones segueix una \fmm{N(100,14)}. Quantes persones tenen un CI entre 108 i 116? I major de 120? \fmpau{PAU C. Valenciana 2001}

\fmm{P(108\leq X\leq116)=P\!\left(\frac{108-100}{14}\leq Z\leq\frac{116-100}{14}\right)=P(0{,}57\leq Z\leq1{,}14)=0{,}8729-0{,}7157=0{,}1572}

\fmm{1200\cdot0{,}1572\approx188{,}6\to} unes 189 persones.

\fmm{P(X\gt120)=P(Z\gt1{,}43)=1-0{,}9236=0{,}0764\Rightarrow1200\cdot0{,}0764\approx92} persones.
\item El temps que tarden els alumnes a fer un examen segueix una \fmm{N(60,10)} (minuts). Quant de temps cal donar perquè l'acabe el 96 \% dels alumnes? \fmpau{PAU País Basc 2020}

Busquem \fmm{t} amb \fmm{P(X\leq t)=0{,}96}. A la taula, el més pròxim a 0,96 és 0,9599 (\fmm{z=1{,}75}).

\fmm{\frac{t-60}{10}=1{,}75\Rightarrow t=60+1{,}75\cdot10=77{,}5} minuts.
\item La nota d'un examen segueix una normal de mitjana 5,8. El 35 \% dels alumnes en trau més de 6,8. Calcula la desviació típica. \fmpau{PAU País Basc 2024}

\fmm{P(X\gt6{,}8)=0{,}35\Rightarrow P\!\left(Z\leq\frac{1}{\sigma}\right)=0{,}65}. A la taula, 0,6517 (\fmm{z=0{,}39}) és el més pròxim.

\fmm{\frac{1}{\sigma}=0{,}39\Rightarrow\sigma=\frac{1}{0{,}39}\approx2{,}56} punts.
\end{enumerate}
\end{fmex}

\Needspace{5.7cm}
\section{Aproximació de la distribució binomial a la normal}\hypertarget{sec9}{}
\begin{fmprop}{Teorema de De Moivre-Laplace}
Quan \fmm{n} és gran, calcular la binomial amb la fórmula és molt pesat (\fmm{P(X\geq40)} amb \fmm{n=100} tindria 61 termes!). Si \fmm{np\geq5} i \fmm{nq\geq5}, la binomial s'assembla molt a una normal amb la mateixa mitjana i desviació típica:

\fmdisp{X\sim B(n,p)\ \approx\ X'\sim N\!\left(np,\ \sqrt{npq}\right)}

\textbf{Correcció de continuïtat}: cada valor enter \fmm{k} de la binomial és una barra que va de \fmm{k-0{,}5} a \fmm{k+0{,}5}. Per això:

\begin{center}\renewcommand{\arraystretch}{1.5}\setlength{\parskip}{0pt}\small\begin{tabularx}{\linewidth}{>{\raggedright\arraybackslash}p{5.2cm}>{\raggedright\arraybackslash}X}\toprule
\textbf{Binomial} & \textbf{Normal}\\\midrule
$P(X=k)$ & $P(k-0{,}5\leq X'\leq k+0{,}5)$\\[3pt]\end{tabularx}\par\nointerlineskip\vskip0pt
\small\begin{tabularx}{\linewidth}{>{\raggedright\arraybackslash}p{5.2cm}>{\raggedright\arraybackslash}X}$P(X\leq k)$ & $P(X'\leq k+0{,}5)$\\[3pt]\end{tabularx}\par\nointerlineskip\vskip0pt
\small\begin{tabularx}{\linewidth}{>{\raggedright\arraybackslash}p{5.2cm}>{\raggedright\arraybackslash}X}$P(X\lt k)=P(X\leq k-1)$ & $P(X'\leq k-0{,}5)$\\[3pt]\end{tabularx}\par\nointerlineskip\vskip0pt
\small\begin{tabularx}{\linewidth}{>{\raggedright\arraybackslash}p{5.2cm}>{\raggedright\arraybackslash}X}$P(X\geq k)$ & $P(X'\geq k-0{,}5)$\\[3pt]\end{tabularx}\par\nointerlineskip\vskip0pt
\small\begin{tabularx}{\linewidth}{>{\raggedright\arraybackslash}p{5.2cm}>{\raggedright\arraybackslash}X}$P(X\gt k)=P(X\geq k+1)$ & $P(X'\geq k+0{,}5)$\\[3pt]\end{tabularx}\par\nointerlineskip\vskip0pt
\small\begin{tabularx}{\linewidth}{>{\raggedright\arraybackslash}p{5.2cm}>{\raggedright\arraybackslash}X}$P(a\leq X\leq b)$ & $P(a-0{,}5\leq X'\leq b+0{,}5)$\\[3pt]\bottomrule\end{tabularx}\end{center}

\begin{fmnote}Truc per a no equivocar-se: escriu primer els enters que \textbf{sí} que entren (per exemple, «més de 499 i com a molt 580» són 500, \dots{}, 580) i després allarga mig punt cap a cada costat: \fmm{499{,}5\leq X'\leq580{,}5}.\end{fmnote}
\end{fmprop}

\Needspace{5\baselineskip}
\begin{fmex}{Exemples resolts}
\begin{enumerate}[label=\textbf{\arabic*.},leftmargin=*]
\item Un creador de vídeos talentós aconsegueix que un vídeo es faça viral amb probabilitat 0,5. Probabilitat que, dels pròxims 100 vídeos, almenys 40 es facen virals. \fmpau{PAU C. Valenciana 2026}

\fmm{X\sim B(100;\,0{,}5)}: \fmm{np=nq=50\geq5}. \fmm{\mu=50}, \fmm{\sigma=\sqrt{100\cdot0{,}5\cdot0{,}5}=5\Rightarrow X'\sim N(50,5)}.

\fmm{P(X\geq40)\approx P(X'\geq39{,}5)=P\!\left(Z\geq\frac{39{,}5-50}{5}\right)=P(Z\geq-2{,}1)=P(Z\leq2{,}1)=0{,}9821}
\item Una asseguradora té 2500 cotxes assegurats i cada un té accident en un any amb probabilitat 0,1. Probabilitat que almenys 270 tinguen accident. \fmpau{PAU Canàries 2022}

\fmm{X\sim B(2500;\,0{,}1)\approx N(250,\ \sqrt{225}=15)}

\fmm{P(X\geq270)\approx P(X'\geq269{,}5)=P(Z\geq1{,}3)=1-0{,}9032=0{,}0968}
\end{enumerate}
\end{fmex}

  
\clearpage
\hypertarget{annex}{}\phantomsection
{\color{fmnavy}\Large\bfseries\colorbox{fmdef}{\color{white}\strut\ Annex\ }\hspace{0.6em}Taules de la binomial i de la normal\par}\vspace{2pt}{\color{fmdef}\titlerule[1.2pt]}
\vspace{6pt}
\label{tbin}{\color{fmgreen}\large\bfseries Taula de la distribució binomial}\par\smallskip
Dona \fmm{P[X=x]=\binom nx\,p^x(1-p)^{n-x}} per a \fmm{n=2,\dots,9}.\newline Si \fmm{p\gt0{,}5}, compta els fracassos: \fmm{Y=n-X\sim B(n,1-p)} (apartat 10.4).\par\smallskip
\begin{center}\linespread{1}\selectfont\scriptsize\setlength{\tabcolsep}{4.2pt}\renewcommand{\arraystretch}{1.0}
\begin{tabular}{cc|ccccccccccccc}\toprule
\multirow{2}{*}{$n$} & \multirow{2}{*}{$x$} & \multicolumn{13}{c}{$p$}\\
 & & \textbf{0,01} & \textbf{0,05} & \textbf{0,10} & \textbf{0,15} & \textbf{0,20} & \textbf{0,25} & \textbf{0,30} & \textbf{1/3} & \textbf{0,35} & \textbf{0,40} & \textbf{0,45} & \textbf{0,49} & \textbf{0,50}\\
\midrule
\multirow{3}{*}{\textbf{2}} & \textbf{0} & 0,9801 & 0,9025 & 0,8100 & 0,7225 & 0,6400 & 0,5625 & 0,4900 & 0,4444 & 0,4225 & 0,3600 & 0,3025 & 0,2601 & 0,2500\\
 & \textbf{1} & 0,0198 & 0,0950 & 0,1800 & 0,2550 & 0,3200 & 0,3750 & 0,4200 & 0,4444 & 0,4550 & 0,4800 & 0,4950 & 0,4998 & 0,5000\\
 & \textbf{2} & 0,0001 & 0,0025 & 0,0100 & 0,0225 & 0,0400 & 0,0625 & 0,0900 & 0,1111 & 0,1225 & 0,1600 & 0,2025 & 0,2401 & 0,2500\\
\midrule
\multirow{4}{*}{\textbf{3}} & \textbf{0} & 0,9703 & 0,8574 & 0,7290 & 0,6141 & 0,5120 & 0,4219 & 0,3430 & 0,2963 & 0,2746 & 0,2160 & 0,1664 & 0,1327 & 0,1250\\
 & \textbf{1} & 0,0294 & 0,1354 & 0,2430 & 0,3251 & 0,3840 & 0,4219 & 0,4410 & 0,4444 & 0,4436 & 0,4320 & 0,4084 & 0,3823 & 0,3750\\
 & \textbf{2} & 0,0003 & 0,0071 & 0,0270 & 0,0574 & 0,0960 & 0,1406 & 0,1890 & 0,2222 & 0,2389 & 0,2880 & 0,3341 & 0,3674 & 0,3750\\
 & \textbf{3} & 0,0000 & 0,0001 & 0,0010 & 0,0034 & 0,0080 & 0,0156 & 0,0270 & 0,0370 & 0,0429 & 0,0640 & 0,0911 & 0,1176 & 0,1250\\
\midrule
\multirow{5}{*}{\textbf{4}} & \textbf{0} & 0,9606 & 0,8145 & 0,6561 & 0,5220 & 0,4096 & 0,3164 & 0,2401 & 0,1975 & 0,1785 & 0,1296 & 0,0915 & 0,0677 & 0,0625\\
 & \textbf{1} & 0,0388 & 0,1715 & 0,2916 & 0,3685 & 0,4096 & 0,4219 & 0,4116 & 0,3951 & 0,3845 & 0,3456 & 0,2995 & 0,2600 & 0,2500\\
 & \textbf{2} & 0,0006 & 0,0135 & 0,0486 & 0,0975 & 0,1536 & 0,2109 & 0,2646 & 0,2963 & 0,3105 & 0,3456 & 0,3675 & 0,3747 & 0,3750\\
 & \textbf{3} & 0,0000 & 0,0005 & 0,0036 & 0,0115 & 0,0256 & 0,0469 & 0,0756 & 0,0988 & 0,1115 & 0,1536 & 0,2005 & 0,2400 & 0,2500\\
 & \textbf{4} & 0,0000 & 0,0000 & 0,0001 & 0,0005 & 0,0016 & 0,0039 & 0,0081 & 0,0123 & 0,0150 & 0,0256 & 0,0410 & 0,0576 & 0,0625\\
\midrule
\multirow{6}{*}{\textbf{5}} & \textbf{0} & 0,9510 & 0,7738 & 0,5905 & 0,4437 & 0,3277 & 0,2373 & 0,1681 & 0,1317 & 0,1160 & 0,0778 & 0,0503 & 0,0345 & 0,0312\\
 & \textbf{1} & 0,0480 & 0,2036 & 0,3281 & 0,3915 & 0,4096 & 0,3955 & 0,3601 & 0,3292 & 0,3124 & 0,2592 & 0,2059 & 0,1657 & 0,1562\\
 & \textbf{2} & 0,0010 & 0,0214 & 0,0729 & 0,1382 & 0,2048 & 0,2637 & 0,3087 & 0,3292 & 0,3364 & 0,3456 & 0,3369 & 0,3185 & 0,3125\\
 & \textbf{3} & 0,0000 & 0,0011 & 0,0081 & 0,0244 & 0,0512 & 0,0879 & 0,1323 & 0,1646 & 0,1811 & 0,2304 & 0,2757 & 0,3060 & 0,3125\\
 & \textbf{4} & 0,0000 & 0,0000 & 0,0005 & 0,0022 & 0,0064 & 0,0146 & 0,0283 & 0,0412 & 0,0488 & 0,0768 & 0,1128 & 0,1470 & 0,1562\\
 & \textbf{5} & 0,0000 & 0,0000 & 0,0000 & 0,0001 & 0,0003 & 0,0010 & 0,0024 & 0,0041 & 0,0053 & 0,0102 & 0,0185 & 0,0282 & 0,0312\\
\midrule
\multirow{7}{*}{\textbf{6}} & \textbf{0} & 0,9415 & 0,7351 & 0,5314 & 0,3771 & 0,2621 & 0,1780 & 0,1176 & 0,0878 & 0,0754 & 0,0467 & 0,0277 & 0,0176 & 0,0156\\
 & \textbf{1} & 0,0571 & 0,2321 & 0,3543 & 0,3993 & 0,3932 & 0,3560 & 0,3025 & 0,2634 & 0,2437 & 0,1866 & 0,1359 & 0,1014 & 0,0938\\
 & \textbf{2} & 0,0014 & 0,0305 & 0,0984 & 0,1762 & 0,2458 & 0,2966 & 0,3241 & 0,3292 & 0,3280 & 0,3110 & 0,2780 & 0,2436 & 0,2344\\
 & \textbf{3} & 0,0000 & 0,0021 & 0,0146 & 0,0415 & 0,0819 & 0,1318 & 0,1852 & 0,2195 & 0,2355 & 0,2765 & 0,3032 & 0,3121 & 0,3125\\
 & \textbf{4} & 0,0000 & 0,0001 & 0,0012 & 0,0055 & 0,0154 & 0,0330 & 0,0595 & 0,0823 & 0,0951 & 0,1382 & 0,1861 & 0,2249 & 0,2344\\
 & \textbf{5} & 0,0000 & 0,0000 & 0,0001 & 0,0004 & 0,0015 & 0,0044 & 0,0102 & 0,0165 & 0,0205 & 0,0369 & 0,0609 & 0,0864 & 0,0938\\
 & \textbf{6} & 0,0000 & 0,0000 & 0,0000 & 0,0000 & 0,0001 & 0,0002 & 0,0007 & 0,0014 & 0,0018 & 0,0041 & 0,0083 & 0,0138 & 0,0156\\
\midrule
\multirow{8}{*}{\textbf{7}} & \textbf{0} & 0,9321 & 0,6983 & 0,4783 & 0,3206 & 0,2097 & 0,1335 & 0,0824 & 0,0585 & 0,0490 & 0,0280 & 0,0152 & 0,0090 & 0,0078\\
 & \textbf{1} & 0,0659 & 0,2573 & 0,3720 & 0,3960 & 0,3670 & 0,3115 & 0,2471 & 0,2048 & 0,1848 & 0,1306 & 0,0872 & 0,0604 & 0,0547\\
 & \textbf{2} & 0,0020 & 0,0406 & 0,1240 & 0,2097 & 0,2753 & 0,3115 & 0,3177 & 0,3073 & 0,2985 & 0,2613 & 0,2140 & 0,1740 & 0,1641\\
 & \textbf{3} & 0,0000 & 0,0036 & 0,0230 & 0,0617 & 0,1147 & 0,1730 & 0,2269 & 0,2561 & 0,2679 & 0,2903 & 0,2918 & 0,2786 & 0,2734\\
 & \textbf{4} & 0,0000 & 0,0002 & 0,0026 & 0,0109 & 0,0287 & 0,0577 & 0,0972 & 0,1280 & 0,1442 & 0,1935 & 0,2388 & 0,2676 & 0,2734\\
 & \textbf{5} & 0,0000 & 0,0000 & 0,0002 & 0,0012 & 0,0043 & 0,0115 & 0,0250 & 0,0384 & 0,0466 & 0,0774 & 0,1172 & 0,1543 & 0,1641\\
 & \textbf{6} & 0,0000 & 0,0000 & 0,0000 & 0,0001 & 0,0004 & 0,0013 & 0,0036 & 0,0064 & 0,0084 & 0,0172 & 0,0320 & 0,0494 & 0,0547\\
 & \textbf{7} & 0,0000 & 0,0000 & 0,0000 & 0,0000 & 0,0000 & 0,0001 & 0,0002 & 0,0005 & 0,0006 & 0,0016 & 0,0037 & 0,0068 & 0,0078\\
\midrule
\multirow{9}{*}{\textbf{8}} & \textbf{0} & 0,9227 & 0,6634 & 0,4305 & 0,2725 & 0,1678 & 0,1001 & 0,0576 & 0,0390 & 0,0319 & 0,0168 & 0,0084 & 0,0046 & 0,0039\\
 & \textbf{1} & 0,0746 & 0,2793 & 0,3826 & 0,3847 & 0,3355 & 0,2670 & 0,1977 & 0,1561 & 0,1373 & 0,0896 & 0,0548 & 0,0352 & 0,0312\\
 & \textbf{2} & 0,0026 & 0,0515 & 0,1488 & 0,2376 & 0,2936 & 0,3115 & 0,2965 & 0,2731 & 0,2587 & 0,2090 & 0,1569 & 0,1183 & 0,1094\\
 & \textbf{3} & 0,0001 & 0,0054 & 0,0331 & 0,0839 & 0,1468 & 0,2076 & 0,2541 & 0,2731 & 0,2786 & 0,2787 & 0,2568 & 0,2273 & 0,2188\\
 & \textbf{4} & 0,0000 & 0,0004 & 0,0046 & 0,0185 & 0,0459 & 0,0865 & 0,1361 & 0,1707 & 0,1875 & 0,2322 & 0,2627 & 0,2730 & 0,2734\\
 & \textbf{5} & 0,0000 & 0,0000 & 0,0004 & 0,0026 & 0,0092 & 0,0231 & 0,0467 & 0,0683 & 0,0808 & 0,1239 & 0,1719 & 0,2098 & 0,2188\\
 & \textbf{6} & 0,0000 & 0,0000 & 0,0000 & 0,0002 & 0,0011 & 0,0038 & 0,0100 & 0,0171 & 0,0217 & 0,0413 & 0,0703 & 0,1008 & 0,1094\\
 & \textbf{7} & 0,0000 & 0,0000 & 0,0000 & 0,0000 & 0,0001 & 0,0004 & 0,0012 & 0,0024 & 0,0033 & 0,0079 & 0,0164 & 0,0277 & 0,0312\\
 & \textbf{8} & 0,0000 & 0,0000 & 0,0000 & 0,0000 & 0,0000 & 0,0000 & 0,0001 & 0,0002 & 0,0002 & 0,0007 & 0,0017 & 0,0033 & 0,0039\\
\midrule
\multirow{10}{*}{\textbf{9}} & \textbf{0} & 0,9135 & 0,6302 & 0,3874 & 0,2316 & 0,1342 & 0,0751 & 0,0404 & 0,0260 & 0,0207 & 0,0101 & 0,0046 & 0,0023 & 0,0020\\
 & \textbf{1} & 0,0830 & 0,2985 & 0,3874 & 0,3679 & 0,3020 & 0,2253 & 0,1556 & 0,1171 & 0,1004 & 0,0605 & 0,0339 & 0,0202 & 0,0176\\
 & \textbf{2} & 0,0034 & 0,0629 & 0,1722 & 0,2597 & 0,3020 & 0,3003 & 0,2668 & 0,2341 & 0,2162 & 0,1612 & 0,1110 & 0,0776 & 0,0703\\
 & \textbf{3} & 0,0001 & 0,0077 & 0,0446 & 0,1069 & 0,1762 & 0,2336 & 0,2668 & 0,2731 & 0,2716 & 0,2508 & 0,2119 & 0,1739 & 0,1641\\
 & \textbf{4} & 0,0000 & 0,0006 & 0,0074 & 0,0283 & 0,0661 & 0,1168 & 0,1715 & 0,2048 & 0,2194 & 0,2508 & 0,2600 & 0,2506 & 0,2461\\
 & \textbf{5} & 0,0000 & 0,0000 & 0,0008 & 0,0050 & 0,0165 & 0,0389 & 0,0735 & 0,1024 & 0,1181 & 0,1672 & 0,2128 & 0,2408 & 0,2461\\
 & \textbf{6} & 0,0000 & 0,0000 & 0,0001 & 0,0006 & 0,0028 & 0,0087 & 0,0210 & 0,0341 & 0,0424 & 0,0743 & 0,1160 & 0,1542 & 0,1641\\
 & \textbf{7} & 0,0000 & 0,0000 & 0,0000 & 0,0000 & 0,0003 & 0,0012 & 0,0039 & 0,0073 & 0,0098 & 0,0212 & 0,0407 & 0,0635 & 0,0703\\
 & \textbf{8} & 0,0000 & 0,0000 & 0,0000 & 0,0000 & 0,0000 & 0,0001 & 0,0004 & 0,0009 & 0,0013 & 0,0035 & 0,0083 & 0,0153 & 0,0176\\
 & \textbf{9} & 0,0000 & 0,0000 & 0,0000 & 0,0000 & 0,0000 & 0,0000 & 0,0000 & 0,0001 & 0,0001 & 0,0003 & 0,0008 & 0,0016 & 0,0020\\
\bottomrule\end{tabular}
\end{center}
\clearpage
\label{tnorm}{\color{fmgreen}\large\bfseries Taula de la distribució normal $N(0,1)$}\par\smallskip
Dona \fmm{F(z)=P(Z\leq z)}, l'àrea a l'esquerra de \fmm{z}. La fila dona les unitats i les dècimes de \fmm{z}; la columna, les centèsimes (apartat 10.7).\par\smallskip
\begin{center}\linespread{1}\selectfont\small\setlength{\tabcolsep}{6pt}\renewcommand{\arraystretch}{1.08}
\begin{tabular}{c|cccccccccc}\toprule
$z$ & \textbf{0,00} & \textbf{0,01} & \textbf{0,02} & \textbf{0,03} & \textbf{0,04} & \textbf{0,05} & \textbf{0,06} & \textbf{0,07} & \textbf{0,08} & \textbf{0,09}\\\midrule
\textbf{0,0} & 0,5000 & 0,5040 & 0,5080 & 0,5120 & 0,5160 & 0,5199 & 0,5239 & 0,5279 & 0,5319 & 0,5359\\
\textbf{0,1} & 0,5398 & 0,5438 & 0,5478 & 0,5517 & 0,5557 & 0,5596 & 0,5636 & 0,5675 & 0,5714 & 0,5753\\
\textbf{0,2} & 0,5793 & 0,5832 & 0,5871 & 0,5910 & 0,5948 & 0,5987 & 0,6026 & 0,6064 & 0,6103 & 0,6141\\
\textbf{0,3} & 0,6179 & 0,6217 & 0,6255 & 0,6293 & 0,6331 & 0,6368 & 0,6406 & 0,6443 & 0,6480 & 0,6517\\
\textbf{0,4} & 0,6554 & 0,6591 & 0,6628 & 0,6664 & 0,6700 & 0,6736 & 0,6772 & 0,6808 & 0,6844 & 0,6879\\[2pt]
\textbf{0,5} & 0,6915 & 0,6950 & 0,6985 & 0,7019 & 0,7054 & 0,7088 & 0,7123 & 0,7157 & 0,7190 & 0,7224\\
\textbf{0,6} & 0,7257 & 0,7291 & 0,7324 & 0,7357 & 0,7389 & 0,7422 & 0,7454 & 0,7486 & 0,7517 & 0,7549\\
\textbf{0,7} & 0,7580 & 0,7611 & 0,7642 & 0,7673 & 0,7704 & 0,7734 & 0,7764 & 0,7794 & 0,7823 & 0,7852\\
\textbf{0,8} & 0,7881 & 0,7910 & 0,7939 & 0,7967 & 0,7995 & 0,8023 & 0,8051 & 0,8078 & 0,8106 & 0,8133\\
\textbf{0,9} & 0,8159 & 0,8186 & 0,8212 & 0,8238 & 0,8264 & 0,8289 & 0,8315 & 0,8340 & 0,8365 & 0,8389\\[2pt]
\textbf{1,0} & 0,8413 & 0,8438 & 0,8461 & 0,8485 & 0,8508 & 0,8531 & 0,8554 & 0,8577 & 0,8599 & 0,8621\\
\textbf{1,1} & 0,8643 & 0,8665 & 0,8686 & 0,8708 & 0,8729 & 0,8749 & 0,8770 & 0,8790 & 0,8810 & 0,8830\\
\textbf{1,2} & 0,8849 & 0,8869 & 0,8888 & 0,8907 & 0,8925 & 0,8944 & 0,8962 & 0,8980 & 0,8997 & 0,9015\\
\textbf{1,3} & 0,9032 & 0,9049 & 0,9066 & 0,9082 & 0,9099 & 0,9115 & 0,9131 & 0,9147 & 0,9162 & 0,9177\\
\textbf{1,4} & 0,9192 & 0,9207 & 0,9222 & 0,9236 & 0,9251 & 0,9265 & 0,9279 & 0,9292 & 0,9306 & 0,9319\\[2pt]
\textbf{1,5} & 0,9332 & 0,9345 & 0,9357 & 0,9370 & 0,9382 & 0,9394 & 0,9406 & 0,9418 & 0,9429 & 0,9441\\
\textbf{1,6} & 0,9452 & 0,9463 & 0,9474 & 0,9484 & 0,9495 & 0,9505 & 0,9515 & 0,9525 & 0,9535 & 0,9545\\
\textbf{1,7} & 0,9554 & 0,9564 & 0,9573 & 0,9582 & 0,9591 & 0,9599 & 0,9608 & 0,9616 & 0,9625 & 0,9633\\
\textbf{1,8} & 0,9641 & 0,9649 & 0,9656 & 0,9664 & 0,9671 & 0,9678 & 0,9686 & 0,9693 & 0,9699 & 0,9706\\
\textbf{1,9} & 0,9713 & 0,9719 & 0,9726 & 0,9732 & 0,9738 & 0,9744 & 0,9750 & 0,9756 & 0,9761 & 0,9767\\[2pt]
\textbf{2,0} & 0,9772 & 0,9778 & 0,9783 & 0,9788 & 0,9793 & 0,9798 & 0,9803 & 0,9808 & 0,9812 & 0,9817\\
\textbf{2,1} & 0,9821 & 0,9826 & 0,9830 & 0,9834 & 0,9838 & 0,9842 & 0,9846 & 0,9850 & 0,9854 & 0,9857\\
\textbf{2,2} & 0,9861 & 0,9864 & 0,9868 & 0,9871 & 0,9875 & 0,9878 & 0,9881 & 0,9884 & 0,9887 & 0,9890\\
\textbf{2,3} & 0,9893 & 0,9896 & 0,9898 & 0,9901 & 0,9904 & 0,9906 & 0,9909 & 0,9911 & 0,9913 & 0,9916\\
\textbf{2,4} & 0,9918 & 0,9920 & 0,9922 & 0,9925 & 0,9927 & 0,9929 & 0,9931 & 0,9932 & 0,9934 & 0,9936\\[2pt]
\textbf{2,5} & 0,9938 & 0,9940 & 0,9941 & 0,9943 & 0,9945 & 0,9946 & 0,9948 & 0,9949 & 0,9951 & 0,9952\\
\textbf{2,6} & 0,9953 & 0,9955 & 0,9956 & 0,9957 & 0,9959 & 0,9960 & 0,9961 & 0,9962 & 0,9963 & 0,9964\\
\textbf{2,7} & 0,9965 & 0,9966 & 0,9967 & 0,9968 & 0,9969 & 0,9970 & 0,9971 & 0,9972 & 0,9973 & 0,9974\\
\textbf{2,8} & 0,9974 & 0,9975 & 0,9976 & 0,9977 & 0,9977 & 0,9978 & 0,9979 & 0,9979 & 0,9980 & 0,9981\\
\textbf{2,9} & 0,9981 & 0,9982 & 0,9982 & 0,9983 & 0,9984 & 0,9984 & 0,9985 & 0,9985 & 0,9986 & 0,9986\\[2pt]
\textbf{3,0} & 0,9987 & 0,9987 & 0,9987 & 0,9988 & 0,9988 & 0,9989 & 0,9989 & 0,9989 & 0,9990 & 0,9990\\
\textbf{3,1} & 0,9990 & 0,9991 & 0,9991 & 0,9991 & 0,9992 & 0,9992 & 0,9992 & 0,9992 & 0,9993 & 0,9993\\
\textbf{3,2} & 0,9993 & 0,9993 & 0,9994 & 0,9994 & 0,9994 & 0,9994 & 0,9994 & 0,9995 & 0,9995 & 0,9995\\
\textbf{3,3} & 0,9995 & 0,9995 & 0,9995 & 0,9996 & 0,9996 & 0,9996 & 0,9996 & 0,9996 & 0,9996 & 0,9997\\
\textbf{3,4} & 0,9997 & 0,9997 & 0,9997 & 0,9997 & 0,9997 & 0,9997 & 0,9997 & 0,9997 & 0,9997 & 0,9998\\[2pt]
\textbf{3,5} & 0,9998 & 0,9998 & 0,9998 & 0,9998 & 0,9998 & 0,9998 & 0,9998 & 0,9998 & 0,9998 & 0,9998\\
\textbf{3,6} & 0,9998 & 0,9998 & 0,9999 & 0,9999 & 0,9999 & 0,9999 & 0,9999 & 0,9999 & 0,9999 & 0,9999\\
\textbf{3,7} & 0,9999 & 0,9999 & 0,9999 & 0,9999 & 0,9999 & 0,9999 & 0,9999 & 0,9999 & 0,9999 & 0,9999\\
\textbf{3,8} & 0,9999 & 0,9999 & 0,9999 & 0,9999 & 0,9999 & 0,9999 & 0,9999 & 0,9999 & 0,9999 & 0,9999\\
\textbf{3,9} & 1,0000 & 1,0000 & 1,0000 & 1,0000 & 1,0000 & 1,0000 & 1,0000 & 1,0000 & 1,0000 & 1,0000\\
\bottomrule\end{tabular}
\end{center}

\end{document}
